%This is my super simple Real Analysis Homework template

\documentclass{article}
\usepackage[utf8]{inputenc}
\usepackage[english]{babel}
\usepackage[]{amsthm} %lets us use \begin{proof}
\usepackage{amsmath}
\usepackage[]{amssymb} %gives us the character \varnothing

\title{Homework 1}
\author{Aine Zhang}
\date\today
%This information doesn't actually show up on your document unless you use the maketitle command below

\begin{document}
\maketitle %This command prints the title based on information entered above

%Section and subsection automatically number unless you put the asterisk next to them.
\section*{Exercise 5.2.8, Problem 1}
\subsection*{i} Let $ f (x) = x^3$ on $[0, 2]$. Compute $\int^2_0f$.
\begin{proof}
$\int^2_0 f = \frac{1}{4}x^4 \vert^{2}_{0} = 4$.
\end{proof}
\subsection*{ii} 
Let
$$
f(x) = 
\begin{cases}
    \frac{1}{q} &  \text{when } x = \frac{p}{q} \text{ in lowest terms, } x \neq 0\\
    0 & \text{otherwise}
\end{cases}
$$
Show that $f$ is integrable, and compute $\int^1_0 f$.
\begin{proof}
Apply problem 3 to 1, we know that here, essentially we can construct a sequence of functions where
$$
f_n  = 
\begin{cases}
    \frac{1}{q} &  \text{when } x = \frac{p}{q},  q \leq n\\
    0 & \text{otherwise}
\end{cases}
$$
$\forall x \text{ and }\forall \epsilon >0, \exists N>0 \text{ such that } |f(x)-f_n(x)| < \epsilon,  \forall n >N$, so we get  
$$ |f(x)-f_n(x)|= 
\begin{cases}
    \frac{1}{q} &  x = \frac{p}{q}, q > n\\
    0 & x = \frac{p}{q}, q \leq n \\
    0 & x \notin \mathbb{Q}
\end{cases}$$
This is true because we defined $f_n$ to be $\frac{1}{q} \text{ for } q>n$. This is all less than $\frac{1}{n}$, which means the function $f$ is integrable, and we know $\int^1_0f = \lim_{k \rightarrow \infty} \int^1_0 f_k =0$.

\end{proof}

\clearpage %Gives us a page break before the next section. Optional.
\section*{Problem 2}
Find another counter-example to show that even the integrability part of Theorem 5.2.10 is false without the word “uniformly”. That is, find a sequence of integrable functions $f_k : [a, b] \rightarrow \mathbb{R}$ which converge pointwise to $f : [a, b] \rightarrow \mathbb{R}$ but f is not Riemann integrable. 
\begin{proof}
Set up a function $f_n$ such that
$$
f_n \rightarrow f(x) = 
\begin{cases}
    1 &  x \in \mathbb{Q}\\
    0 & \text{otherwise}
\end{cases}
$$
$f_n$ is constructed as:
$$
f_n(x_n)= 
\begin{cases}
    1 &  \{ x_1, x_2...x_n\} \in \mathbb{Q}\\
    0 & \text{otherwise}
\end{cases}
$$
Since $\mathbb{Q} \cap [0, 1]$ is countably finite by Cantor's diagonal argument, we can set $\forall x, \lim_{n \rightarrow + \infty}f_n(x) = f(x)$. If $x \in \mathbb{Q}, x = x_n$, by definition $f_n(x_n) = 1$.  $\forall n \geq N$, which gives us $f_n(x_n) \text{ as } n \rightarrow \infty $ is eventually $1$, which gives $\lim_{n \rightarrow \infty} f_n(x) = 1 = f(x).$ Thus, we know $f_n$ converges to $f$ pointwise, but not uniformly. Each $f_n$ itself is Riemann integrable because the property $\mathbb{Q} \cap [0, 1]$ is countably finite means there are only finite discontinuities. 

\end{proof}
\clearpage
\section*{Problem 3}
Let $f : [a, b] \rightarrow \mathbb{R}$ be an integrable function. Let $\{ x_n \}^{\infty}_{n=1}=1 \subset (a, b)$ be a sequence of distinct numbers, and $c_n \in \mathbb{R}$ be a sequence of numbers with $ \lim_{n \rightarrow \infty} c_n = 0 $. Define
$$
g(x) = 
\begin{cases}
    f(x) + c_n & x = x_n \text{ for some } n\\
    f(x) & \text{otherwise}
\end{cases}
$$
Prove $g$ is integrable. (Hint: Use Theorem 5.2.10)
\begin{proof}
Given $g(x)$, we can construct a function where 
$$
f_n(x_n) = 
\begin{cases}
    g(x_1) = f(x) + c_1 & x = x_1 \\
    g(x_2) = f(x) + c_2 & x = x_1 \\
    ...\\
    f(x) & \text{otherwise}
\end{cases}
$$, which is just
$$
f_n(x_n) = 
\begin{cases}
    g(x_n) = f(x) + c_n & x = x_n \text{ for some } n \\
    f(x) & \text{otherwise}
\end{cases}
$$
Then, we know $ = $
$$
|f(x)-f_n(x_n) | = 
\begin{cases}
    c_n & \text{for some } n \\
    0 & \text{otherwise}
\end{cases}
$$
Since the problem gives us $ \lim_{n \rightarrow \infty} c_n = 0 $, we know $\lim_{n \rightarrow \infty}|f(x)-f_n(x_n) | \leq \epsilon \text{ for some } \epsilon \geq 0$. This means the function $f_n = g$ converges pointwise to $f(x)$, so by theorem 5.2.10, the function is Riemann integrable.
\end{proof}
\clearpage
\section*{Problem 4}
\subsection*{i}
Let $f : [a, b] \rightarrow \mathbb{R}$ be a bounded function and let $\mathcal{P}$ be a partition of $[a, b]$. Let $\mathcal{P}'$ be a refinement of $\mathcal{P}$. Show that $L(f, \mathcal{P}) \leq L(f, \mathcal{P}') \leq U (f, \mathcal{P}') \leq U (f, \mathcal{P})$. \\
\begin{proof}
    We can use induction to show the statement is true. Let us say we insert a point $b$ into the existing partition $\mathcal{P}$ so that the original partition $\mathcal{P} = \{ a_0, a_1...a_i, a_j...a_n \}$ becomes a refinement $\mathcal{P} = \{ a_0, a_1...a_i, b, a_j...a_n \}$. \\

    %We want to show $L(f, \mathcal{P}) \leq L(f, \mathcal{P}') $
    By the definition of a lower sum, we know $$L(f, \mathcal{P}) = \sum^n_{i=1}\inf_{x \in [a_{i-1}, a_i]}f(x)(a_i-a_{i-1}).$$ Let us abbreviate $\inf_{x \in [a_{i-1}, a_i]}f(x)(a_i-a_{i-1}) $ as $\inf_{x \in [a_{i-1}, a_i]}$ here. By expanding the definition of the lower sum, we know it is essentially $$ (a_1-a_0)\inf_{x \in [a_0, a_1]} + (a_2-a_1) \inf_{x \in [a_{1}, a_2]} + ... + (a_j - a_i) \inf_{x \in [a_{i}, a_j]} +...(a_n - a_{n-1})\inf_{x \in [a_{n-1}, a_n]},$$
    and by the same definition applied to the refinement with $b$, we get  
    $$L(f, \mathcal{P}') = \sum^n_{i=1}\inf_{x \in [a_{i-1}, a_i]}f(x)(a_i-a_{i-1}), $$ 
    which will expand to $$ (a_1-a_0)\inf_{x \in [a_0, a_1]} + (a_2-a_1) \inf_{x \in [a_{1}, a_2]} + ... + (b - a_i) \inf_{x \in [a_i, b]}+(a_j - b) \inf_{x \in [b, a_j]} +...(a_n - a_{n-1})\inf_{x \in [a_{n-1}, a_n]}.$$ \\
    Since $[a_i, b]$ and $[b, a_j]$ will both be smaller intervals than $[a_i, a_j]$, we can expect $\inf_{x \in [a_{i}, b]}f(x) \leq \inf_{x \in [a_{i}, a_j]}f(x) \text{  and} \inf_{x \in [b, a_j]}f(x) \leq \inf_{x \in [a_{i}, a_j]}f(x) ,$ as well as $(b - a_i) \leq (a_j-a_i)$ and $(a_j - b) \leq (a_j-a_i)$. So this gives us  $L(f, \mathcal{P}) \leq L(f, \mathcal{P}_{\text{with inserted b}}') $. By induction, any refinement can be thought of inserting multiple such $b$s into the original partition $\mathcal{P}$. So we get the general application $L(f, \mathcal{P}) \leq L(f, \mathcal{P}') $
\end{proof}
\subsection*{ii}
Show that for any two partitions $\mathcal{P}$ and $\mathcal{P}'$, $L(f, \mathcal{P}) \leq U(f, \mathcal{P}')$. \\
\begin{proof}
    Let's construct a partition $\mathcal{P''}$ where $\mathcal{P}'' = \mathcal{P } \cup \mathcal{P'}$, then $\mathcal{P}''$ is a refinement of both $\mathcal{P}$ and $\mathcal{P}'$. By the previous question, we get $L(f, \mathcal{P}) \leq L(f, \mathcal{P}'')$, $U(f, \mathcal{P}'')\leq U (f, \mathcal{P}')$ and by definition of lower and upper sums, $L (f, \mathcal{P}'') \leq U (f, \mathcal{P}'')$. By chaining these together, we get $L(f, \mathcal{P}) \leq L(f, \mathcal{P}'') \leq U (f, \mathcal{P}'') \leq U (f, \mathcal{P}')$, so this gives us $L(f, \mathcal{P}) \leq U (f, \mathcal{P}')$ for any two partitions $\mathcal{P}$ and $\mathcal{P'}.$
\end{proof}


\end{document}